\section{Introduction}\label{sec:introduction}

The small spectrum of a degenerating metric records both the topology of
the degeneration and the energy required to join its limiting harmonic
fields. When several local scales tend to zero independently, an estimate
at the largest small eigenvalue can lose the smaller modes. We prove a
relative comparison for the complete small two-form spectrum of actual
Ricci-flat metrics on a compact Calabi--Yau conifold family. The family
has eight nodes, two lost divisor directions, and six independent
vanishing-period coordinates. It permits arbitrary separation of the
two positive spectral scales.

The central threefold and its small resolution are fixed throughout.
They arise from the anticanonical contraction of
\begin{equation}\label{eq:ambient-intro}
 Z=\PP_{\PP^1\times\PP^1}(\OO\oplus\OO(1,-1))
       \times_{\PP^1\times\PP^1}
       \PP_{\PP^1\times\PP^1}(\OO\oplus\OO(1,-1)).
\end{equation}
For a general anticanonical hypersurface $Y\subset Z$, chosen also
nondegenerate with respect to its toric polytope, the induced morphism
$q:Y\to X_0$ contracts eight disjoint $(-1,-1)$ curves. They form two
groups $I_f,I_g$ of four curves of the same respective homology types.
This compact construction and its Picard data are established in
\cite[Appendix~A]{wuebben-chambers}. We retain the exact polarization
$\cL_0=(-K_{\bar Z})|_{X_0}$, where $\bar Z$ is the anticanonical
model. Section~\ref{sec:geometry} describes the contraction and proves
the required multiparameter smoothing statement.

There is a ball $B$ about zero in
\begin{equation}\label{eq:period-base-intro}
 R=\left\{p\in\CC^8:
       \sum_{i\in I_f}p_i=\sum_{i\in I_g}p_i=0\right\}
\end{equation}
and a polarized compact family $\pi:\mathcal X\to B$ with central
fibre $X_0$. A nowhere-zero relative holomorphic three-form $\Omega$
has oriented local vanishing periods $p_i$. In node coordinates the
equations are
\begin{equation}\label{eq:node-unit-intro}
 z_1^2+z_2^2+z_3^2+z_4^2=\tau_i(p),\qquad
 p_i=\tau_i(p)u_i(p),
\end{equation}
where the $u_i$ are holomorphic and bounded above and away from zero.
The smooth locus is exactly $B^\circ=\{p\in B:\prod_i p_i\ne0\}$.
Thus the node parameters are controlled by their periods even when one
is smaller than every prescribed power of another. A first-order
description of the deformation space would not give this conclusion.

Let $g_p$ be the Ricci-flat K\"ahler metric in $c_1(\cL_p)$ on
$X_p$, for $p\in B^\circ$. We use the real Hodge Laplacian
$\Delta_p=\dd\delta_p+\delta_p\dd$, with its usual closed-manifold
operator domain. Define
\begin{equation}\label{eq:finite-pair-intro}
 \kappa_i(p)=\frac{2\pi}{|\log|\tau_i(p)||},\qquad
 w_\alpha(p)=\sum_{i\in I_\alpha}\kappa_i(p),\qquad
 E_p=\begin{pmatrix}w_f&0\\0&w_g\end{pmatrix},\qquad
 M=\frac45\begin{pmatrix}9&-1\\-1&9\end{pmatrix}.
\end{equation}
The energy is the capacity-weighted squared restriction of the lost
divisor classes to the links. The mass is the central harmonic mass
after orthogonal projection off all surviving cohomology. Both are
constructed independently of the actual positive small eigenvalues.

To specify the comparison of eigenspaces, fix a compact central exterior
$C_*$ by deleting node balls of a small positive radius. Its central
metric defines the $L^2$ projection $\Pi_*$ onto distributionally closed
two-forms, with unrestricted tangential trace. Let $V=H^2(Y;\RR)$,
let $G$ be its central harmonic mass, and let $K\subset V$ be the
two-dimensional subspace of classes surviving on the smooth fibres.
Put $W=K^{\perp_G}$. The cohomology of the exterior identifies its
closed projection with a class in $V$. The observation map is
\begin{equation}\label{eq:observer-intro}
 \cA u=\pr_W^G[\Pi_*(u|_{C_*})].
\end{equation}
Here and below the restrictions to $C_*$ use fixed smooth exterior
identifications. The closed projection and its cohomology map are
bounded on exterior $L^2$ forms. Formula~\eqref{eq:observer-intro}
therefore has a precise meaning for form-domain inputs; it does not
take arbitrary cycle integrals of an unclosed $L^2$ form. We use the
$G$-orthogonal lifts of the quotient classes $f,g$ as coordinates on $W$.

\begin{theorem}\label{thm:main}
After shrinking $B$, there are $\gamma>0$ and a function $e(p)\to0$
as $p\to0$ in $B^\circ$ with the following properties.

The harmonic two-form space of $(X_p,g_p)$ has dimension two. Its
cohomological identification with $K$ transports its mass to a form
converging to $G|_K$. Exactly two positive two-form eigenvalues lie below
$\gamma$. If $\mu_1(p)\le\mu_2(p)$ are the generalized eigenvalues
of $(E_p,M)$, these eigenvalues satisfy
\begin{equation}\label{eq:main-rates}
 (1-e(p))\mu_j(p)\le\lambda_{2,j}^+(p)
                 \le(1+e(p))\mu_j(p),\qquad j=1,2.
\end{equation}
The first positive one-form eigenvalue and the third positive two-form
eigenvalue are at least $\gamma$.

The two-dimensional positive small spectral space $S_p$ consists of
coexact primitive real $(1,1)$ forms. The restriction $\cA:S_p\to W$
is an isomorphism. Its transported mass and energy,
\begin{equation}\label{eq:observed-pair-intro}
 N_p(x,y)=\langle\cA^{-1}x,\cA^{-1}y\rangle_{g_p},\qquad
 Q_p(x,y)=\langle\dd\cA^{-1}x,\dd\cA^{-1}y\rangle_{g_p},
\end{equation}
satisfy the quadratic-form inequalities
\begin{equation}\label{eq:main-forms}
 (1-e(p))M\le N_p\le(1+e(p))M,\qquad
 (1-e(p))E_p\le Q_p\le(1+e(p))E_p.
\end{equation}
All conclusions hold throughout $B^\circ$ as $p\to0$, with no bound
on the ratios of the nonzero node parameters or of the positive model
eigenvalues.
\end{theorem}

The whole-form conclusion retains mixed pairings in the prescribed
coordinates. If $u_1,u_2$ are actual orthonormal eigenforms, $C_p$ has
columns $\cA u_j$, and $\Lambda_p=\operatorname{diag}
(\lambda_{2,1}^+,\lambda_{2,2}^+)$, it gives
\begin{equation}\label{eq:marked-matrix-intro}
 \|C_p^TMC_p-I\|\le e(p),\qquad
 \|\Lambda_p^{-1/2}(C_p^TE_pC_p-\Lambda_p)
                          \Lambda_p^{-1/2}\|\le e(p),
\end{equation}
after changing the error function by a fixed factor. In particular, an
off-diagonal energy error is bounded by a vanishing multiple of
$\sqrt{\lambda_{2,1}^+\lambda_{2,2}^+}$. No explicit rate is claimed
for the compactness-derived part of $e(p)$.

The finite eigenvalues are
\begin{equation}\label{eq:finite-roots-intro}
 \mu_\pm=
 \frac{9(w_f+w_g)\pm\sqrt{81(w_f-w_g)^2+4w_fw_g}}{128}.
\end{equation}
They satisfy $0<\mu_-/\mu_+\le4/5$. The corresponding actual
eigenvalues are consequently simple for all sufficiently small
parameters. A relative subspace estimate identifies each actual
eigenline, through $\cA$, with its finite generalized eigenline, with
error tending to zero uniformly in $w_f/w_g$.

For example, the admissible path
\begin{equation}\label{eq:separated-path-intro}
 p(h)=(a,a,a,-3a;b,b,b,-3b),\qquad
 a=e^{-h^{-1}},\quad b=e^{-h^{-2}},\quad h\downarrow0,
\end{equation}
has actual asymptotics
\begin{equation}\label{eq:separated-spectrum-intro}
 \lambda_{2,1}^+(p(h))\sim\frac{10\pi}{9}h^2,
 \qquad
 \lambda_{2,2}^+(p(h))\sim\frac{9\pi}{8}h.
\end{equation}
More generally $a=e^{-F(h)}$, $b=e^{-G(h)}$, with positive
$F,G\to\infty$ and $F/G\to0$, gives coefficients
$9\pi/(8F)$ and $10\pi/(9G)$. These are paths in a holomorphic
multiparameter family; they need not extend as holomorphic arcs at zero.

\subsection{Relative reduction and earlier results}

Relative small-spectrum comparison has substantial precedents. Burger's
graph theorem for degenerating hyperbolic surfaces already permits
unequal short-geodesic lengths and gives multiplicative scalar
eigenvalue asymptotics \cite[Theorem~1.1]{burger}. In cusp surgery,
Albin--Rochon--Sher identify positive small Hodge modes by a restriction
difference and compute their leading rescaled energy
\cite[Section~3, Lemma~3.2 and Theorem~3.3]{ars-cusps}.
Anné--Takahashi's conic analysis retains the cap and critical extended
solutions needed for a complete spectral count
\cite[Theorems~A--B]{anne-takahashi}. The geometry and parameter domains
of these theorems differ from the polarized Ricci-flat smoothing used
here.

For actual Ricci-flat conifolds, Hein--Sun establish the central conical
metric asymptotics \cite[Theorem~1.4]{hein-sun}. Benabida proves
polyhomogeneous smoothing expansions along prescribed rays
\cite[Theorem~C]{benabida}. Pacini's scalar gluing estimates already
allow independent scales on distinct cap components
\cite[Theorems~11.7 and 12.3]{pacini}. We verify the scalar decomposition,
initial K\"ahler family, global inverse and nonlinear correction on the
whole period domain, then prove the complete two-form comparison there.
The fixed-ray spectrum in \cite{wuebben-hodge} supplies the central
geometric normalization, not the unequal-scale theorem by substitution.

The metric proof first identifies all fifteen scalar harmonic modes
of growth below three on a Stenzel cap. Nine are pluriharmonic; the
remaining six are derivatives of complex orthogonal automorphisms.
A global scalar inverse retains these modes separately from a
positive-weight remainder. In the nonlinear equation we use the
finite automorphisms themselves: their exact preservation of the
Stenzel volume removes the nondecaying quadratic terms. This gives
the actual Ricci-flat comparison on every cap and intervening annulus
without a smallest-scale loss.

For the two-form problem, corrected logarithmic trials have exactly
the capacity energy, including zero contribution from a link on
which the period vanishes. A weighted estimate on each independent
cap prevents low-energy mass concentration in degrees one and two.
A fixed exterior closed projection then gives a lower bound in the
input's own energy norm. Multiplicative metric comparison transfers
the full count and gaps to the actual metrics, while the exterior
projection preserves the cohomological marking.

The abstract argument in Section~\ref{sec:relative} isolates three
geometric estimates: linear trials, a weighted lower energy inequality
for the specified observation, and capture of low-energy mass. It
allows the exact kernel to change and obtains relative energy on the
entire observed low space. The final step uses a determinant on the
quotient by the common kernel. It is useful precisely when an additive
mass-plus-energy comparison would lose the weakest direction. The
geometric estimates, rather than the finite matrix factorization alone,
are the analytic content of an application.

The hierarchy of the finite forms retains the capacities, the full
restriction map and the harmonic mass. Its relation to the classical
attachment sequence is explicit. It is distinct from the Leray
filtration of a fixed adiabatic fibration
\cite{mazzeo-melrose}, and from the corrected foliation spectral-sequence
statements of Álvarez López--Kordyukov \cite{alvarez-kordyukov}.
No new nearby-cycle functor or canonical integral lattice on a general
small eigenspace is asserted.
